
En
effet, soient \(\eta\) le point g\'en\'erique de \(S\) et \(K\) la cl\^oture alg\'ebrique de \(k(\eta)\). Il est clair que \(\Spec(K)\) est la limite du syst\`eme projectif filtrant \((S_\alpha)\) des sch\'emas affines int\`egres munis d'un morphisme fini et fid\`element plat \(S_\alpha\to U_\alpha\) o\`u \(U_\alpha\) est un ouvert non-vide de \(S\). Le normalis\'e de \(X_K\) provient donc d'un \(S_\alpha\)-sch\'ema \(X'\) pour un \(\alpha\), le morphisme fini et surjectif \(X'_K\to X_K\) d'un \(S'_\alpha\)-morphisme fini et surjectif \(X'\to X_{S_\alpha}\) (EGA IV 8.8.2, 8.10.5). Enfin, comme \(X'_K\) est normal et int\`egre, il existe un ouvert \(S'\) de \(S_\alpha\) tel que \(X'_{S'}\) soit normal et int\`egre sur \(S'\) (EGA IV 9.9.5).

\textbf{Remarque 3.10.} Conservant les notations et conditions de 3.9, supposons qu'il existe un morphisme r\'esolvant \(X'_K\to X_K\) (EGA IV 7.9.1), (par exemple si \(\dim(X_K)\le2\), d'apr\`es Abhyankhar). Alors un raisonnement similaire montre qu'il existe un morphisme fini et fid\`element plat \(S'\to U\), un sch\'ema \(X'\) lisse sur \(S'\) et un morphisme propre surjectif birationnel \(X'\to X_{S'}\).

Dans la d\'emonstration de la premi\`ere partie de 3.8, on s'est ramen\'e au cas o\`u \(S\) est int\`egre et \(Y\) est normal et int\`egre sur \(S\), donc \(Y\) est int\`egre. Si \(\varphi=0\), alors \(X=Y\) et l'assertion est triviale. Sinon, quitte \`a restreindre \(S\), on peut supposer que pour tout \(s\in S\), on a \(\varphi_s\ne0\), donc \(\varphi_s\) est \({O}_{Y_s}\)-r\'egulier et la dimension \(r\) de \(Y_s\) est \(\ge3\).

Soit alors \(\mathcal L\) une famille limit\'ee de classes de faisceaux inversibles sur les fibres de \(X/S\). Soit \(M_K\) un faisceau inversible sur une fibre g\'eom\'etrique \(Y_K\) de \(Y/S\), avec le polyn\^ome de Hilbert \(\chi(M_K(n))=\sum_{i=0}^{r}a_i\binom{n+i}{i}\), alors le polyn\^ome de Hilbert de \(f^*M_K\) est \(\sum_{i=0}^{r-1}a_{i+1}\binom{n+i}{i}\). Lorsque la classe de \(M_K\) parcourt \((f^*)^{-1}(\mathcal L)\), c-\`a-d la
